\documentclass[11pt,a4paper]{article}
% Shared preamble for the English probability and statistics documents.
% Based on the original course preamble. Compile with pdfLaTeX.
\usepackage[top=1.5cm,bottom=2cm,left=2cm,right=2cm]{geometry}
\usepackage[T1]{fontenc}
\usepackage{lmodern}
\usepackage[utf8]{inputenc}
\usepackage[english]{babel}

\usepackage{amsmath}
\usepackage{amssymb}
\usepackage{amsthm}
\usepackage{mathrsfs}
\usepackage{mathtools}
\usepackage{etoolbox}
\usepackage{tikz}
\usetikzlibrary{shapes.multipart,arrows.meta}
\usepackage{xcolor}
\usepackage{enumerate}
\usepackage[shortlabels]{enumitem}
\usepackage{hyperref}
\usepackage{version}
\usepackage{bbold}
\usepackage{fontawesome5}
\usepackage{fancyhdr}
\lfoot{}
\usepackage{natbib}
\usepackage{graphicx}
\usepackage{pdfpages}
\usepackage{authblk}
    
\usepackage{blindtext}
\usepackage{multicol}
\usepackage{collcell}
\usepackage{wrapfig}
\usepackage{blkarray}
\usepackage{cancel}
\usepackage{appendix}
\usepackage{float}
\makeatletter
\input{tex/assets/slashbox.sty}
\makeatother
\hypersetup{unicode,colorlinks,citecolor=red,filecolor=black,
  linkcolor=red,urlcolor=cyan}
\setlength{\columnsep}{1cm}
\pagenumbering{gobble}

\def\numtext#1{%
  \ifcase#1\or one\or two\or three\or four\or five\or six\or seven\or eight\or nine\or ten\or
  eleven\or twelve\or thirteen\or fourteen\or fifteen\or sixteen\or seventeen\or eighteen\or nineteen\or twenty\or Lots\fi}

\newcommand{\Byte}[1]{%
  \def\nodecontents{}%
  \foreach[count=\l] \k in {#1} {\xappto{\nodecontents}{\noexpand\nodepart{\numtext{\l}}\k}}%
  \node[rectangle split,rectangle split parts=8,rectangle split horizontal,draw] at (2,2){\nodecontents};}

\newcommand{\D}{{\mathrm{d}}}
\newcommand{\cE}{{\mathcal{E}}}
\newcommand{\cG}{{\mathcal{G}}}
\newcommand{\cX}{{\mathcal{X}}}
\newcommand{\cY}{{\mathcal{Y}}}
\newcommand{\cB}{{\mathcal{B}}}
\newcommand{\cF}{{\mathcal{F}}}
\newcommand{\cO}{{\mathcal{O}}}
\newcommand{\cS}{{\mathcal{S}}}
\newcommand{\cT}{{\mathcal{T}}}
\newcommand{\given}{\, \big| \,}
\newcommand{\cL}{{\mathscr{L}}}
\newcommand{\cD}{{\mathcal{D}}}
\newcommand{\cA}{{\mathscr{A}}}
\newcommand{\cP}{{\mathscr{P}}}
\newcommand{\cR}{{\mathscr{R}}}
\newcommand{\A}{{\mathcal A}}
\newcommand{\one}{\mathbb{1}}
\newcommand{\mathbbm}[1]{\mathbb{#1}}

\providecommand{\C}{}
\renewcommand{\C}{\mathbb C}
\newcommand{\R}{\mathbb R}
\newcommand{\Q}{\mathbb Q}
\newcommand{\N}{\mathbb N}
\newcommand{\T}{\mathbb T}
\newcommand{\F}{\mathbb F}
\renewcommand{\P}{\mathbb P}
\newcommand{\Z}{\mathbb Z}
\newcommand{\E}{\mathbb {E}}
\newcommand{\Pbb}[1]{\P\left( #1\right)}
\newcommand{\Ebb}[1]{\E\left[ #1\right]}

\newcommand{\lb}{\left (}
\newcommand{\rb}{\right )}
\newcommand{\lbb}{\left [}
\newcommand{\rbb}{\right ]}
\newcommand{\labs}{\left |}
\newcommand{\rabs}{\right |}
\newcommand{\lbrb}[1]{\lb #1 \rb}
\newcommand{\lbbrbb}[1]{\lbb#1\rbb}
\newcommand{\labsrabs}[1]{\labs#1\rabs}
\newcommand{\lbbrb}[1]{\lbb#1\rb}
\newcommand{\lbrbb}[1]{\lb#1\rbb}
\newcommand{\langlerangle}[1]{\langle#1\rangle}
\newcommand{\lbcurly}{\{}
\newcommand{\rbcurly}{\}}
\newcommand{\curly}[1]{\lbcurly#1\rbcurly}
\newcommand{\abs}{\labsrabs}
\newcommand{\ind}[1]{1_{#1}}
\newcommand{\indep}{\perp \!\!\! \perp}
\renewcommand{\d}{\, \mathrm{d}}

\newtheorem{theorem}{Theorem}
\newtheorem{problemCnt}{Lemma}
\theoremstyle{definition}
\newtheorem{problem}[problemCnt]{Problem}
\newtheorem{problembook}[problemCnt]{\faBook \,  Problem}
\newtheorem{problempencil}[problemCnt]{\faPencil* \,  Problem}
\newtheorem{problemstar}[problemCnt]{\faStar\,  Problem}
\newtheorem{problemstarpencil}[problemCnt]{\faStar \, \faPencil* \,   Problem}
\newtheorem{probleminteresting}[problemCnt]{\faSmile[regular] \,   Problem}
\newtheorem{solution}{Solution}

\newcommand{\fwcell}[1]{\makebox[\arraystretch\normalbaselineskip]{$#1$}}
\newenvironment{enumeratei}{\begin{enumerate}[(i)]}{\end{enumerate}}
\newenvironment{enumeratea}{\begin{enumerate}[(a)]}{\end{enumerate}}
\newenvironment{enumeraten}{\begin{enumerate}[\upshape 1.]}{\end{enumerate}}

% Compatibility helpers needed by a few files in the combined archive.
\AtBeginDocument{\noindent{\footnotesize English translation prepared with ChatGPT (OpenAI, GPT-6), 26 October 2026, for translation of the Bulgarian text into English and LaTeX formatting.
The Bulgarian originals remain the reference; translation may contain errors.}\par\smallskip}
\newcommand{\MissingGraphic}[1]{%
  \begin{center}\fbox{\parbox{0.82\linewidth}{\centering Missing source image: \texttt{#1}}}\end{center}}



\pagestyle{fancy}\fancyhf{}\lhead{Computer Science: Test 1 solutions}\rhead{20 April 2024}
\setlength{\headheight}{15pt}\pagenumbering{arabic}\cfoot{\thepage}
\begin{document}
\noindent\emph{English transcription of the handwritten solution notes. The calculations and original level of detail are retained; apparent source slips are identified in the transcription notes.}

\section*{Problem 1}
\subsection*{1. Conditional probability}
There are two yellow, three green and four blue balls. A yellow ball means rolling one die, green means two dice, and blue means three dice. By Bayes' formula,
\begin{align*}
\P(\text{green}\mid\text{sum }3)
&=\frac{\P(\text{sum }3\mid\text{green})\P(\text{green})}{\P(\text{sum }3)}\\
&=\frac{\frac2{36}\cdot\frac39}
{\frac16\cdot\frac29+\frac2{36}\cdot\frac39+\frac1{216}\cdot\frac49}
=\frac9{28}\approx32.14\%.
\end{align*}
\subsection*{2. Five dice}
\[
\P(\max\ge3)=1-\P(\text{all outcomes are 1 or 2})=1-\left(\frac26\right)^5.
\]
Let $C$ be the event that at least one number is divisible by 5, and $D$ the event that at least one number is divisible by 2. Then
\[
\P(C^c)=\left(\frac56\right)^5,\qquad
\P(D^c)=\left(\frac12\right)^5,\qquad
\P(C^c\cap D^c)=\left(\frac26\right)^5.
\]
We seek
\begin{align*}
\P(C\cap D)&=\P(C)+\P(D)-\P(C\cup D)\\
&=1-\P(C^c)+1-\P(D^c)-\bigl(1-\P(C^c\cap D^c)\bigr)\\
&=1-\left(\frac56\right)^5-\left(\frac12\right)^5
+\left(\frac16\right)^5\approx56.70\%.\tag{Note 1}
\end{align*}
\subsection*{3. Seven balls in seven boxes}
Two boxes contain three balls each and one box contains one ball when $X_3=2$, so
\[
\P(X_3=2)=\frac{7\binom73\,6\binom43\,5/2}{7^7}\approx1.78\%.
\]
When $X_3=1$, the nonzero occupancies are $(3,4)$, $(3,1,1,1,1)$, $(3,2,2)$ or $(3,2,1,1)$. Thus
\begin{align*}
\P(X_3=1)&=\frac{7\binom73\,6}{7^7}
+\frac{7\binom73\,6\cdot5\cdot4\cdot3}{7^7}\\
&\quad+\frac{7\binom73\,6\binom42\,5/2}{7^7}
+\frac{7\binom73\,6\binom42\,5\cdot4}{7^7}
\approx34.99\%.
\end{align*}
The source then writes
\[
\P(X_3=0)=1-\P(X_3=0\text{ or }1)\approx63.23\%.\tag{Note 2}
\]
Consequently,
\[
\E X_3=\P(X_3=1)+2\P(X_3=2)\approx0.3855.
\]
Also,
\[
\sum_{i=1}^7i\E X_i=\E\left[\sum_{i=0}^7iX_i\right]
=\E[\text{number of balls}]=7.
\]
Similarly,
\[
\sum_{i=0}^7\E X_i=\E[X_0+\cdots+X_7]
=\E[\text{number of boxes}]=7.
\]

\section*{Problem 2}
Let $Y$ be the number of correct answers among the questions the student knows, and $Z$ the number of correct answers among those the student does not know. Then
\[
Y\sim\operatorname{Bin}(10,9/10),\qquad
Z\sim\operatorname{Bin}(5,1/5),\qquad X=Y+Z.
\]
\subsection*{1. More than twelve correct answers; moments}
\begin{align*}
\P(X>12)&=\sum_{k=3}^5\P(Z=k)\P(Y>12-k)\\
&=\P(Z=3)\P(Y=10)\\
&\quad+\P(Z=4)[\P(Y=10)+\P(Y=9)]\\
&\quad+\P(Z=5)[\P(Y=10)+\P(Y=9)+\P(Y=8)].
\end{align*}
\[
\E X=\E Y+\E Z=10\cdot\frac9{10}+5\cdot\frac15=10.
\]
By independence,
\[
DX=DY+DZ=10\cdot\frac9{10}\cdot\frac1{10}
+5\cdot\frac15\cdot\frac45=\frac{17}{10}.
\]
\subsection*{2. Conditional probability}
By Bayes' formula,
\begin{align*}
\P(Y=3\mid X=3)
&=\frac{\P(X=3\mid Y=3)\P(Y=3)}
{\sum_k\P(X=3\mid Y=k)\P(Y=k)}\\
&=\frac{\P(Z=0)\P(Y=3)}{\sum_{k=0}^3\P(Z=3-k)\P(Y=k)}\\
&=\frac{(4/5)^5\binom{10}3(9/10)^3(1/10)^7}
{\displaystyle\sum_{k=0}^3\binom5{3-k}(1/5)^{3-k}(4/5)^{k+2}
\binom{10}k(9/10)^k(1/10)^{10-k}}=\cdots.
\end{align*}

\section*{Problem 3}
\subsection*{1. Joint distribution}
\[
\begin{array}{c|ccc|c}
Z\backslash X&0&1&2&\P(Z=z)\\\hline
0&1/9&0&2/9&1/3\\
1&2/9&1/9&0&1/3\\
2&0&2/9&1/9&1/3\\\hline
\P(X=x)&1/3&1/3&1/3&1
\end{array}
\]
The source calls the second variable $Y$ in the following calculation (the variable denoted $Z$ in the table):
\[
\E X=\E Y=1,\qquad
\E[XY]=1\cdot1\cdot\frac19+2\cdot1\cdot\frac29+2\cdot2\cdot\frac19=1.
\]
Hence $\operatorname{Cov}(X,Y)=\E[XY]-\E X\E Y=0$ and $\operatorname{Cor}(X,Y)=0$.

\subsection*{2. Generating function}
\begin{align*}
g_X(s)&=\tfrac16(s+\cdots+s^6),\\
g'_X(s)&=\tfrac16(1+2s+\cdots+6s^5),&g'_X(1)&=\tfrac72,\\
g''_X(s)&=\tfrac16(2+6s+12s^2+20s^3+30s^4),&g''_X(1)&=\tfrac{70}6.
\end{align*}
Thus
\[
\E X=g'_X(1)=\frac72,\qquad
DX=g''_X(1)+g'_X(1)-[g'_X(1)]^2
=\frac{70}6+\frac72-\frac{49}4=\frac{35}{12}.
\]

\subsection*{3. Independence}
If $X,Y$ are independent, then $\P(X=k,Y=m)=\P(X=k)\P(Y=m)$, so
\begin{align*}
g_{X+Y}(s)&=\sum_{k=0}^\infty s^k\P(X+Y=k)\\
&=\sum_{k=0}^\infty s^k\sum_{m=0}^k\P(X=m)\P(Y=k-m)\\
&=\sum_{m=0}^\infty\sum_{k=m}^\infty
s^m\P(X=m)s^{k-m}\P(Y=k-m)\\
&=\sum_{m=0}^\infty s^m\P(X=m)\sum_{i=0}^\infty s^i\P(Y=i)
=g_X(s)g_Y(s),
\end{align*}
where $i=k-m$. The converse is false: $g_{X+Y}=g_Xg_Y$ does not imply independence. The variables $X,Z$ from part 1 give a counterexample. Verify this!

\section*{Problem 4}
Let $p_n(k)=\P(k\in E_n)$, the probability that $k$ is selected when the initial choice is from $\{1,\ldots,n\}$. The key observation is
\[
p_n(k)=\underbrace{p_n(k+1)}_{\substack{\text{select }k\text{ without selecting }k+1\\\text{(why is this probability }p_n(k+1)\text{?)}}}
+\underbrace{p_n(k+1)\frac1k}_{\text{select }k+1\text{ and then }k}.
\]
Thus $p_n(k)=\frac{k+1}{k}p_n(k+1)$. Since $p_n(n)=1/n$,
\[
p_n(k)=\frac{k+1}{k}\frac{k+2}{k+1}\cdots\frac n{n-1}p_n(n)=\frac1k.
\]
\paragraph{2.} Using the preceding observation,
\[
\P(2\in E_n\mid3\notin E_n)=\frac{p_n(3)}{1-p_n(3)}
=\frac{1/3}{2/3}=\frac12.
\]
\paragraph{3.} Let $X_i=\mathbf1_{\{i\text{ was selected}\}}\sim\operatorname{Ber}(1/i)$.
By linearity,
\[
\E|E_n|=\E[X_1+\cdots+X_n]=\sum_{i=1}^n\E X_i
=1+\frac12+\cdots+\frac1n.
\]
\paragraph{Remark 1.} In fact, $X_i$ and $X_j$ are independent for $i\ne j$!
\paragraph{Remark 2.} This procedure is actually the reverse of the Reservoir Sampler (see the exercise classes), which also explains Remark 1.

\section*{Transcription notes}
\begin{enumerate}
\item The source's final term $(1/6)^5$ conflicts with the preceding $(2/6)^5$. The reported $56.70\%$ agrees with $(2/6)^5$.
\item The source writes ``$X_3=0$ or $1$'' in the complement. It should refer to $X_3=1$ or $X_3=2$.
\end{enumerate}
\end{document}
